\subsection{The rational rank of commutative group}

DEFINITION 1. The rational rank of a commutative group G is the dimension of the vector Q-space \(G \otimes_{\mathbf{Z}} \mathbf{Q}\).

This dimension may also be defined as the least upper bound (finite or infinite) of the cardinals r such that there exist r elements of G which are linearly independent over Z (Algebra, Chapter II, § 7, no. 10, Proposition 26). The rational rank of G is zero if and only if G is a torsion group. For a subgroup of an additive group \(R^{n}\) , the notion of rational rank coincides with that defined in General Topology, Chapter VII, § 1.

In the rest of this paragraph, we shall denote by \(r(G)\) the rational rank of the commutative group \(G\) . If \(G'\) is a subgroup of \(G\) , then (since \(Q\) is a flat \(Z\) -module) we have the additive equation

\[
r (\mathbf {G}) = r \left(\mathbf {G} ^ {\prime}\right) + r \left(\mathbf {G} / \mathbf {G} ^ {\prime}\right).\tag{4}
\]

PROPOSITION 3. Let G be a totally ordered commutative group and H a subgroup of G. If \(h(\mathbf{G})\) and \(h(\mathbf{H})\) denote the heights of G and H (§ 4, no. 4), then the inequality

\[
h (\mathbf {G}) \leqslant h (\mathbf {H}) + r (\mathbf {G} / \mathbf {H})\tag{5}
\]

holds.

In fact, let \(G_{0} \subset G_{1} \subset \cdots \subset G_{n}\) be a strictly increasing sequence of isolated subgroups of G. It is necessary to establish the inequality

\[
n \leqslant h (\mathrm{H}) + r (\mathrm{G} / \mathrm{H}).\tag{6}
\]

It is obvious for n = 0. Suppose \(n \geqslant 1\) and let us argue by induction on n.~Applying the induction hypothesis to the group \(G_{n-1}\) and its subgroup \(H \cap G_{n-1}\) , we obtain

\[
n - 1 \leqslant h (\mathrm{H} \cap \mathrm{G} _ {n - 1}) + r (\mathrm{G} _ {n - 1} / (\mathrm{H} \cap \mathrm{G} _ {n - 1})).\tag{7}
\]

Then we distinguish two cases:

\begin{enumerate}
\def\labelenumi{(\alph{enumi})}
\tightlist
\item
  \(\mathbf{H} \cap \mathbf{G}_{n-1} = \mathbf{H}\) , in other words \(\mathbf{H} \subset \mathbf{G}_{n-1}\) . Inequality (7) may be written
\end{enumerate}

\[
n \leqslant h (\mathrm{H}) + r (\mathrm{G} _ {n - 1} / \mathrm{H}) + 1.\tag{8}
\]

Now \(\mathbf{G} / \mathbf{G}_{n - 1}\) is a totally ordered group not reduced to 0; it is therefore not a torsion group and \(r(\mathrm{G} / \mathrm{G}_{n - 1}) \geqslant 1\) . Whence by (4), \(r(\mathrm{G} / \mathrm{H}) \geqslant r(\mathrm{G}_{n - 1} / \mathrm{H}) + 1\) . Substituting in (8), we certainly obtain the desired inequality (6).

\begin{enumerate}
\def\labelenumi{(\alph{enumi})}
\setcounter{enumi}{1}
\tightlist
\item
  \(\mathbf{H} \cap \mathbf{G}_{n-1} \neq \mathbf{H}\) . As \(\mathbf{H} \cap \mathbf{G}_{n-1}\) is an isolated subgroup of \(\mathbf{H}\) , we conclude that \(h(\mathbf{H}) \geqslant h(\mathbf{H} \cap \mathbf{G}_{n-1}) + 1\) . On the other hand, obviously \(r(\mathbf{G}/\mathbf{H}) \geqslant r(\mathbf{G}_{n-1}/(\mathbf{H} \cap \mathbf{G}_{n-1}))\) . Substituting in (7), we again obtain (6).
\end{enumerate}

COROLLARY. For every totally ordered commutative group \(\mathbf{G}\) , \(h(\mathbf{G}) \leqslant r(\mathbf{G})\) .

We set \(\mathbf{H} = \{0\}\) in Proposition 3.

PROPOSITION 4. Let G be a totally ordered commutative group. Suppose that G is finitely generated and that \(h(\mathbf{G}) = r(\mathbf{G})\) . Then G is isomorphic to \(\mathbf{Z}^{r(\mathbf{G})}\) ordered lexicographically.

Let us write \(r = r(\mathbf{G}) = h(\mathbf{G})\) . If \(r = 0\) , then \(\mathbf{G} = \{0\}\) . If \(r = 1\) , the structure of finitely generated commutative groups shows that there is an isomorphism \(j\) of \(\mathbf{G}\) onto \(\mathbf{Z}\) (Algebra, Chapter VII, § 4, no. 6, Theorem 3). Now \(\mathbf{Z}\) has only two total orderings compatible with its group structure, namely the usual ordering and its opposite. Hence \(j\) or \(-j\) is an isomorphism of the ordered group \(\mathbf{G}\) onto \(\mathbf{Z}\) with the usual ordering.

Suppose now that \(r \geqslant 2\) and let us argue by induction on \(r\) . Let \(H\) be an isolated subgroup of \(G\) of height \(r - 1\) . Then \(r(H) + r(G/H) = r\) (formula (4)), \(r(H) \geqslant h(H) = r - 1\) and \(r(G/H) \geqslant h(G/H) = 1\) (Corollary to Proposition 3), whence \(r(H) = r - 1\) and \(r(G/H) = 1\) . The induction hypothesis shows that \(H\) is isomorphic to \(Z^{r-1}\) ordered lexicographically and the case \(r = 1\) shows that \(G/H\) is isomorphic to \(Z\) . As \(Z\) is a free \(Z\) -module, \(H\) is a direct factor of \(G\) (Algebra, Chapter II, § 1, no. 11, Proposition 21). The following lemma then shows that \(G\) is isomorphic (not canonically) to the lexicographical product \(H \times (G/H)\) , which completes the proof.

LEMMA 2. Let H be an isolated subgroup of a totally ordered commutative group G. If H is a direct factor of G, the ordered group G is isomorphic to the group (G/H) × H ordered lexicographically.

Let \(j\) be an isomorphism of \((\mathrm{G} / \mathrm{H}) \times \mathrm{H}\) onto \(\mathbf{G}\) such that \(j(0, x) = x\) for all \(x \in \mathrm{H}\) and \(j(y, x)\) belongs to the coset \(y\) of \(\mathrm{H}\) . As \((\mathrm{G} / \mathrm{H}) \times \mathrm{H}\) is totally ordered, it amounts to showing that \(j\) is increasing (Set Theory, Chapter III, § 1, no. 12, Proposition 11). Let \((y, x)\) be an element \(\geqslant 0\) of \((\mathrm{G} / \mathrm{H}) \times \mathrm{H}\) ordered lexicographically. If \(y > 0\) , the coset of \(\mathrm{H}\) containing \(j(y, x)\) is an element \(> 0\) , whence \(j(y, x) > 0\) , for, otherwise, \(y \leqslant 0\) (§ 4, no. 2, Proposition 3). If \(y = 0\) and \(x \geqslant 0\) , then \(j(y, x) = x \geqslant 0\) . Hence \(j\) is certainly increasing.

